% ═══════════════════════════════════════════════════════════════════════
% Section 16 – Market Clearing Mathematical Models
% (SCUC / SCED / LMP)
%
% Reference: Southern China regional spot-market implementation rules (2025 V1.0)
%            Section 2.6 day-ahead energy-market clearing
% ═══════════════════════════════════════════════════════════════════════
\newpage
\section{Market Clearing Mathematical Models}\label{sec:market-clearing}

\subsection{Problem Form and Runtime Scope}

This chapter documents the day-ahead market-clearing stack built in
\texttt{src/analysis/market\_simulation.cpp}.  The runtime sequence is a layered
optimization pipeline:
\begin{itemize}
  \item SCUC as a multi-period MILP,
  \item SCED as an LP re-dispatch with fixed commitment,
  \item LMP extraction from LP dual variables on the cleared network model.
\end{itemize}

\subsection{Code Map}

The current implementation is concentrated in:
\begin{itemize}
  \item \texttt{src/analysis/market\_simulation.cpp} for SCUC, SCED, LMP model assembly, and reporting,
  \item \texttt{src/engine/adapter\_registry.cpp}, \texttt{src/engine/native\_adapters.cpp}, and \texttt{src/engine/external\_adapters.cpp} for backend routing,
  \item \texttt{src/engine/branch\_and\_cut.cpp} and associated LP backends for native MILP and LP execution when external solvers are unavailable.
\end{itemize}

\subsection{Pseudocode Map}

The detailed algorithm chapter does not yet include a dedicated market-clearing
pseudocode block.  The solver infrastructure underneath this chapter is covered
in Section~\ref{sec:engine-solver-system} and Section~\ref{sec:native-branch-and-cut}.
This chapter therefore remains the primary problem-form and runtime-sequencing
reference for SCUC, SCED, and LMP.

\subsection{Current Status}

The market-clearing route is implemented, but it remains one of the more active
integration surfaces in the repository because it depends on MILP backends,
pricing duals, and network-constraint assembly.  This chapter should therefore
be read as an implementation-facing description of the current solver stack,
not as a claim that every possible market-model extension is already active.

This section presents the complete mathematical formulations for day-ahead
electricity market clearing implemented in \texttt{market\_simulation.cpp}.
The models follow the \emph{Southern China Regional Electricity Market Spot Energy
Trading Implementation Rules} (Southern China regional spot-market rules,
2025~V1.0), Section~2.6.

The day-ahead clearing consists of three sequential stages:
\begin{enumerate}
  \item \textbf{SCUC} (Security-Constrained Unit Commitment, §2.6.3):
        Determines the $T$-period commitment schedule via mixed-integer linear
        programming (MILP).
  \item \textbf{SCED} (Security-Constrained Economic Dispatch, §2.6.4):
        Fixes the commitment from SCUC and re-optimises dispatch as a linear
        program (LP).
  \item \textbf{LMP} (Locational Marginal Price, §2.6.5--2.6.6):
        Computes nodal prices from the LP dual variables of a
        $\delta$-neighbourhood re-dispatch model.
\end{enumerate}

%-------------------------------------------------------------------
\subsection{Nomenclature}\label{sec:market-nomenclature}
%-------------------------------------------------------------------

\subsubsection{Sets and Indices}

\begin{longtable}{lL{10.5cm}}
\toprule
\textbf{Symbol} & \textbf{Description} \\
\midrule
\endfirsthead
\toprule
\textbf{Symbol} & \textbf{Description} \\
\midrule
\endhead
$\mathcal{G}$       & Set of all generators (thermal, nuclear, hydro), indexed by $i$, $|\mathcal{G}|=N$; does not include storage \\
$\mathcal{W}$       & Set of wind farms, indexed by $w$, $|\mathcal{W}|=N_W$ \\
$\mathcal{S}$       & Set of solar PV plants, indexed by $s$, $|\mathcal{S}|=N_{PV}$ \\
$\mathcal{R}$       & $\mathcal{W} \cup \mathcal{S}$, set of all renewable units, $|\mathcal{R}|=N_R$ \\
$\mathcal{E}$       & Set of energy storage units, indexed by $e$, $|\mathcal{E}|=N_{ES}$ \\
$\mathcal{D}$       & Set of loads, indexed by $k$, $|\mathcal{D}|=K$ \\
$\mathcal{L}$       & Set of AC monitored transmission lines, $|\mathcal{L}|=N_L$ \\
$\mathcal{S}_{\mathrm{sec}}$ & Set of monitored sections/corridors, $|\mathcal{S}_{\mathrm{sec}}|=N_S$ \\
$\mathcal{DC}$      & Set of intra-regional DC interconnections, $|\mathcal{DC}|=N_{TD}$ \\
$\mathcal{T}_{O}$   & Set of inter-regional (external) tie-lines, $|\mathcal{T}_{O}|=N_T$ \\
$\mathcal{A}$       & Set of balancing areas (provinces), $|\mathcal{A}|=N_A$ \\
$\mathcal{T}$       & Set of time periods, $|\mathcal{T}|=T$ (D-day: 96 $\times$ 15\,min; D+1: 2 peak/valley; total 98) \\
$\mathcal{M}_i$     & Set of bid segments for generator $i$, $|\mathcal{M}_i|=N_M$ \\
$\mathcal{I}_s$     & Set of must-run / CHP / commissioning units \\
$\mathcal{H}$       & Set of hydro stations with reservoir, $|\mathcal{H}|=N_H$ \\
$\mathcal{J}$       & Set of generator groups, indexed by $j$ \\
\bottomrule
\end{longtable}

\subsubsection{Decision Variables}

\begin{longtable}{lL{10.5cm}}
\toprule
\textbf{Variable} & \textbf{Description} \\
\midrule
\endfirsthead
\toprule
\textbf{Variable} & \textbf{Description} \\
\midrule
\endhead
$P_{i,t}$           & Generation output of unit $i$ in period $t$ [MW] \\
$\Delta P_{i,t,m}$  & Dispatch in bid segment $m$ of unit $i$ in period $t$ [MW] \\
$u_{i,t}$ ($\alpha_{i,t}$) & Commitment status of unit $i$ in period $t$ (binary: 0/1) \\
$v_{i,t}$           & Startup indicator of unit $i$ in period $t$ (binary: 1 iff $u_{i,t}=1, u_{i,t-1}=0$) \\
$w_{i,t}$           & Shutdown indicator of unit $i$ in period $t$ (binary: 1 iff $u_{i,t}=0, u_{i,t-1}=1$) \\
$R^{\mathrm{spin}}_{i,t}$ & Spinning reserve of unit $i$ in period $t$ [MW] \\
$R^{\mathrm{up}}_{i,t}$   & Regulation-up reserve of unit $i$ in period $t$ [MW] \\
$R^{\mathrm{dn}}_{i,t}$   & Regulation-down reserve of unit $i$ in period $t$ [MW] \\
$P^{W}_{w,t}$       & Wind generation of farm $w$ in period $t$ [MW] \\
$P^{PV}_{s,t}$      & Solar PV generation of plant $s$ in period $t$ [MW] \\
$P^{\mathrm{ch}}_{e,t}$  & Storage charging power of unit $e$ in period $t$ [MW] \\
$P^{\mathrm{dis}}_{e,t}$ & Storage discharging power of unit $e$ in period $t$ [MW] \\
$E_{e,t}$           & State of charge (SOC) of storage $e$ at end of period $t$ [MWh] \\
$\xi^{+}_{e,t}$, $\xi^{-}_{e,t}$ & Binary indicators for storage charge/discharge status \\
$T^{DC}_{d,t}$      & Intra-regional DC line $d$ power transfer in period $t$ [MW] \\
$T^{O}_{j,t}$       & Inter-regional tie-line $j$ scheduled power in period $t$ [MW] \\
$SL^{+}_{l,t}$, $SL^{-}_{l,t}$ & Forward/reverse line flow slack for line $l$ in period $t$ \\
$SL^{+}_{s,t}$, $SL^{-}_{s,t}$ & Forward/reverse section flow slack for section $s$ in period $t$ \\
$SL_k$              & Tie-line priority schedule slack for component $k$ \\
$P^{d}_{r,t}$       & Renewable curtailment of unit $r$ in period $t$ [MW] \\
$P^{d,\mathrm{hydro}}_{i,t}$ & Hydro spillage power of station $i$ in period $t$ [MW] \\
$\mathrm{LS}_t$     & Load shedding in period $t$ [MW] \\
$\mathrm{GC}_t$     & Generation curtailment in period $t$ [MW] \\
\bottomrule
\end{longtable}

\subsubsection{Parameters}

\begin{longtable}{lL{10.5cm}}
\toprule
\textbf{Parameter} & \textbf{Description} \\
\midrule
\endfirsthead
\toprule
\textbf{Parameter} & \textbf{Description} \\
\midrule
\endhead
$C_{i,m}$           & Bid price of generator $i$ in segment $m$ [\textyen/MWh] \\
$C^{U}_{i}$         & Startup cost of unit $i$ (hot/warm/cold state dependent) [\textyen] \\
$C^{\min,p}_{i}$    & Minimum-output cost of unit $i$ [\textyen/h]; applies when $u_{i,t}=1$ \\
$P^{\min}_i$        & Minimum stable output of unit $i$ [MW] \\
$P^{\max}_i$        & Maximum capacity of unit $i$ [MW] \\
$\Delta P^{\min}_{i,m}$, $\Delta P^{\max}_{i,m}$ & Bid segment $m$ lower/upper bound [MW] \\
$\Delta P^{U}_i$    & Ramp-up rate of unit $i$ [MW/period] \\
$\Delta P^{D}_i$    & Ramp-down rate of unit $i$ [MW/period] \\
$T^{U}_i$           & Minimum up-time of unit $i$ [periods] \\
$T^{D}_i$           & Minimum down-time of unit $i$ [periods] \\
$N^{\max,\mathrm{start}}_i$ & Maximum startup count of unit $i$ over the horizon \\
$N^{\max,\mathrm{stop}}_i$  & Maximum shutdown count of unit $i$ over the horizon \\
$UD_i$              & Startup process duration (periods from sync to min output) \\
$DD_i$              & Shutdown process duration (periods from min output to desync) \\
$D^{a}_{t}$         & System load in balancing area $a$, period $t$ [MW] \\
$R^{U,a}_{t}$       & Positive reserve requirement for area $a$ in period $t$ [MW] \\
$R^{D,a}_{t}$       & Negative reserve requirement for area $a$ in period $t$ [MW] \\
$R^{\mathrm{red},a}_{t}$ & Network-constrained reserve reduction for area $a$ [MW] \\
$R^{\mathrm{pf}}_t$ & Primary frequency regulation reserve requirement [MW] \\
$P^{\max}_l$, $P^{\min}_l$ & Forward/reverse thermal limits of line $l$ [MW] \\
$P^{\max}_s$, $P^{\min}_s$ & Forward/reverse limits of section $s$ [MW] \\
$G_{l,k}$           & PTDF of bus $k$ on line $l$ \\
$G_{s,k}$           & PTDF of bus $k$ on section $s$ \\
$\eta^{\mathrm{ch}}_e$, $\eta^{\mathrm{dis}}_e$ & Charging/discharging efficiency of storage $e$ \\
$\bar{E}_e$, $\underline{E}_e$ & SOC upper/lower limits of storage $e$ [MWh] \\
$E^{\mathrm{ini}}_e$ & Initial SOC of storage $e$ [MWh] \\
$E^{\mathrm{fin}}_e$ & Required final SOC of storage $e$ [MWh] \\
$N^{\mathrm{cycle}}_e$ & Maximum daily charge/discharge cycles of storage $e$ \\
$P^{\mathrm{st}}_{r,t}$ & Short-term forecast power of renewable $r$ in period $t$ [MW] \\
$\alpha_n$          & Renewable minimum output coefficient (province-specific) \\
$\lambda^{\mathrm{dis}}_e$, $\lambda^{\mathrm{ch}}_e$ & Storage discharge/charge bid prices [\textyen/MWh] \\
$P_{\mathrm{gwf}}$  & Inter-provincial transmission fee (incl.\ network losses, wheeling) [\textyen/MWh] \\
$M_1$               & Network flow constraint relaxation penalty [\textyen/MW] \\
$M_2$               & Renewable curtailment penalty [\textyen/MW] \\
$M_3$               & Hydro spillage penalty [\textyen/MW] \\
$M_4$               & Tie-line priority schedule relaxation penalty [\textyen/MW] \\
$\Delta t$          & Period length [hours] (0.25 for 15-min periods) \\
$\delta$            & LMP $\delta$-neighbourhood parameter (fraction) \\
\bottomrule
\end{longtable}

%===================================================================
\subsection{SCUC -- Security-Constrained Unit Commitment (§2.6.3)}
\label{sec:scuc}
%===================================================================

\subsubsection{Objective Function}\label{sec:scuc-obj}

The SCUC minimises total system cost over the planning horizon $T$
(§2.6.3 Objective):
\begin{equation}\label{eq:scuc-obj}
\boxed{
\begin{aligned}
\min \;\; & \sum_{i=1}^{N}\sum_{t=1}^{T}
  \bigl[ C_i(P_{i,t}) + C^{U}_{i,t} + C^{\min,p}_{i,t} \bigr]
  \;+\; \sum_{i=1}^{n}\sum_{t=1}^{T} P_{\mathrm{gwf}} \cdot P_{L,i,t} \\
  &+\; M_1 \sum_{l=1}^{N_L}\sum_{t=1}^{T}(SL^{+}_{l,t}+SL^{-}_{l,t})
  \;+\; M_1 \sum_{s=1}^{N_S}\sum_{t=1}^{T}(SL^{+}_{s,t}+SL^{-}_{s,t}) \\
  &+\; M_2 \sum_{r=1}^{N_R}\sum_{t=1}^{T} P^{d}_{r,t}
  \;+\; M_3 \sum_{i=1}^{N_H}\sum_{t=1}^{T} P^{d,\mathrm{hydro}}_{i,t}
  \;+\; M_4 \sum_{k=1}^{a}\sum_{t=1}^{T} SL_k \\
  &+\; \sum_{e=1}^{N_{ES}}\sum_{t=1}^{T}
    \bigl(\lambda^{\mathrm{dis}}_e \cdot P^{\mathrm{dis}}_{e,t}
    + \lambda^{\mathrm{ch}}_e \cdot P^{\mathrm{ch}}_{e,t}\bigr)
\end{aligned}
}
\end{equation}

\paragraph{Generator operating cost (§2.6.3.7).}
The operating cost is a piecewise-linear function of the bid segments:
\begin{equation}\label{eq:gen-cost}
  C_i(P_{i,t}) = \sum_{m=1}^{N_M} C_{i,m} \cdot \Delta P_{i,t,m}
\end{equation}

\paragraph{Startup cost (§2.6.3.13).}
Startup cost depends on off-time duration (cold/warm/hot state):
\begin{equation}\label{eq:startup-cost}
  C^{U}_{i,t} = C^{U}_i \cdot v_{i,t}
\end{equation}

\paragraph{Minimum-output cost.}
Minimum-output cost applies when the unit is committed:
\begin{equation}\label{eq:noload-cost}
  C^{\min,p}_{i,t} = C^{\min,p}_i \cdot u_{i,t} \cdot \Delta t
\end{equation}

\paragraph{C++ mapping.}
In \texttt{build\_scuc\_formulation()}:
\begin{itemize}[nosep]
  \item Segment cost: \texttt{lp.c(seg\_idx) = price * period\_length\_hr}
  \item Startup cost: \texttt{lp.c(su\_idx) += startup\_cost\_hot}
  \item Min-output cost: \texttt{lp.c(ig\_idx) += no\_load\_cost / intervals\_per\_hour}
  \item $M_1$ slack: \texttt{lp.c(SL\_LINE/SEC) = M1\_line\_slack\_penalty}
  \item Load shedding: \texttt{lp.c(LOAD\_SHED) = VOLL * period\_length\_hr}
  \item Gen curtailment: \texttt{lp.c(GEN\_CURT) = VOCC * period\_length\_hr}
  \item Storage bids: \texttt{lp.c(PSTO\_OUT/IN)} stores the discharge or charge bid coefficient \(\lambda_e^{\mathrm{dis}}\) or \(\lambda_e^{\mathrm{ch}}\)
\end{itemize}

%-------------------------------------------------------------------
\subsubsection{Constraint 1: Power Balance (§2.6.3.1)}\label{sec:scuc-balance}

For each balancing area $a$ and period $t$:
\begin{equation}\label{eq:power-balance}
  \sum_{i \in \mathcal{G}^a} P_{i,t}
  + \sum_{j=1}^{NT^a} T^{O,a}_{j,t}
  + \sum_{k=1}^{NTI^a} T^{DC,a}_{k,t}
  = D^{a}_{t}
\end{equation}
where $P_{i,t}$ includes generators, renewables, and storage net output
(discharge minus charge). The implementation adds load-shedding
($\mathrm{LS}_t$) and generation-curtailment ($\mathrm{GC}_t$) slack
variables to the balance.

%-------------------------------------------------------------------
\subsubsection{Constraint 2: Positive Reserve (§2.6.3.2)}\label{sec:scuc-pos-reserve}

For each province $a$ and period $t$:
\begin{equation}\label{eq:pos-reserve}
  \sum_{i \in \mathcal{G}^a}
    u_{i,t} \cdot P^{\max,a}_{i,t}
  + \sum_{j=1}^{NT^a} T^{O,a}_{j,t}
  + \sum_{k=1}^{NTI^a} T^{DC,a}_{k,t}
  \;\geq\;
  D^{a}_{t} + R^{U,a}_{t} - R^{\mathrm{red},a}_{t}
\end{equation}

In C++, reserves are modelled via explicit reserve variables with coupling:
\begin{equation}\label{eq:reserve-coupling}
  P_{i,t} + R^{\mathrm{spin}}_{i,t} + R^{\mathrm{up}}_{i,t}
  \;\leq\; P^{\max}_i \cdot u_{i,t}
  \qquad \forall\, i \in \mathcal{G},\; t \in \mathcal{T}
\end{equation}

%-------------------------------------------------------------------
\subsubsection{Constraint 3: Negative Reserve (§2.6.3.3)}\label{sec:scuc-neg-reserve}

For each province $a$ and period $t$:
\begin{equation}\label{eq:neg-reserve}
  \sum_{i \in \mathcal{G}^a}
    \bigl(P_{i,t} - \beta_{i,t} \cdot P^{\min,a}_{i,t}\bigr)
  + \sum_{q} \beta_{q,t} \cdot P^{\max,a}_{q,t}
  + \sum_{j} T^{O,a}_{j,t}
  + \sum_{k} T^{DC,a}_{k,t}
  \;\leq\;
  D^{a}_{t} - R^{D,a}_{t}
\end{equation}
where $\beta_{i,t}$ indicates whether unit $i$ can provide negative reserve
(thermal: same as $u_{i,t}$; hydro/VPP: separately determined).

%-------------------------------------------------------------------
\subsubsection{Constraint 4: Primary Frequency Regulation Reserve (§2.6.3.4)}\label{sec:scuc-pf-reserve}

\begin{equation}\label{eq:pf-reserve}
  \sum_{f=1}^{NF} P^{\mathrm{pf}}_{f,t}
  + \sum_{p=1}^{NHP} P^{\mathrm{pf}}_{p,t}
  + \sum_{h=1}^{NHC} P^{\mathrm{pf}}_{h,t}
  \;\geq\; R^{\mathrm{pf}}_t
  \qquad \forall\, t
\end{equation}
where $NF$, $NHP$, $NHC$ are the counts of thermal, pumped-storage,
and conventional hydro units, and each unit's PFR capacity is:
\begin{equation}
  P^{\mathrm{pf}}_{f,t} = \min\bigl(\alpha^{\mathrm{pf}}_f \cdot P^{\max}_{f,t},\;
  P^{\max}_{f,t} - P_{f,t}\bigr)
\end{equation}

%-------------------------------------------------------------------
\subsubsection{Constraint 5: Must-Run Units (§2.6.3.5)}\label{sec:scuc-mustrun}

\begin{equation}\label{eq:must-run}
  u_{i,t} = 1 \qquad \forall\, i \in \mathcal{I}_s,\; t \in \mathcal{T}
\end{equation}
$\mathcal{I}_s$ includes must-run generators, CHP units (during heating
periods), and commissioning units.

%-------------------------------------------------------------------
\subsubsection{Constraint 6: Generator Output Expression (§2.6.3.6)}\label{sec:scuc-gen-output}

Output is decomposed into a minimum block and piecewise segments:
\begin{equation}\label{eq:gen-output}
  P_{i,t} = P^{\min}_i \cdot u_{i,t} + \sum_{m=1}^{N_M} \Delta P_{i,t,m}
\end{equation}
\begin{equation}\label{eq:seg-bounds}
  \Delta P^{\min}_{i,m} \;\leq\; \Delta P_{i,t,m} \;\leq\; \Delta P^{\max}_{i,m}
  \qquad \forall\, m \in \mathcal{M}_i
\end{equation}

%-------------------------------------------------------------------
\subsubsection{Constraint 7: Generator Output Limits (§2.6.3.8)}\label{sec:scuc-gen-limits}

\begin{equation}\label{eq:gen-limits}
  P^{\min}_{i,t} \cdot u_{i,t}
  \;\leq\; P_{i,t} \;\leq\;
  P^{\max}_{i,t} \cdot u_{i,t}
  \qquad \forall\, i \in \mathcal{G},\; t \in \mathcal{T}
\end{equation}
For units in startup/shutdown transition (§2.6.3.8 detailed modelling),
the limits account for $UD_i$ and $DD_i$ ramp trajectories:
\begin{equation}\label{eq:startup-traj}
  P^{\min}_{i,t} =
  \begin{cases}
    P^{\min}_i & \text{if } u_{i,t}=1 \text{ (steady state)} \\
    \text{startup trajectory value} & \text{if startup in progress}
  \end{cases}
\end{equation}

%-------------------------------------------------------------------
\subsubsection{Constraint 8: Generator Group Output Limits (§2.6.3.9)}\label{sec:scuc-group-output}

\begin{equation}\label{eq:group-output}
  P^{\min}_{j,t} \;\leq\;
  \sum_{i \in j} P_{i,t}
  \;\leq\; P^{\max}_{j,t}
  \qquad \forall\, j \in \mathcal{J},\; t \in \mathcal{T}
\end{equation}

%-------------------------------------------------------------------
\subsubsection{Constraint 9: Generator Group Energy Limits (§2.6.3.10)}\label{sec:scuc-group-energy}

\begin{equation}\label{eq:group-energy}
  Q^{\min}_j \;\leq\;
  \sum_{i \in j} \sum_{t=1}^{T} Q_{i,t}
  \;\leq\; Q^{\max}_j
\end{equation}
where $Q_{i,t} = P_{i,t} \cdot \Delta t$ is the cleared energy of unit $i$
in period $t$.

%-------------------------------------------------------------------
\subsubsection{Constraint 10: Ramp Rate (§2.6.3.11)}\label{sec:scuc-ramp}

\begin{equation}\label{eq:ramp-up}
  P_{i,t} - P_{i,t-1} \;\leq\; \Delta P^{U}_i + P^{\max}_i \cdot v_{i,t}
  \qquad \forall\, i,\; t \geq 2
\end{equation}
\begin{equation}\label{eq:ramp-down}
  P_{i,t-1} - P_{i,t} \;\leq\; \Delta P^{D}_i + P^{\max}_i \cdot w_{i,t}
  \qquad \forall\, i,\; t \geq 2
\end{equation}
The startup/shutdown terms allow larger ramps during transitions.

%-------------------------------------------------------------------
\subsubsection{Constraint 11: Minimum Up/Down Time (§2.6.3.12)}\label{sec:scuc-minud}

Commitment transition logic:
\begin{equation}\label{eq:transition}
  u_{i,t} - u_{i,t-1} = v_{i,t} - w_{i,t}
  \qquad \forall\, i,\; t \geq 2
\end{equation}

Minimum up-time: if a unit started up, it must stay on for $T^{U}_i$ periods:
\begin{equation}\label{eq:min-up}
  \sum_{\tau=\max(1,\,t-T^{U}_i+1)}^{t} v_{i,\tau} \;\leq\; u_{i,t}
  \qquad \forall\, i,\; t
\end{equation}

Minimum down-time: if a unit shut down, it must stay off for $T^{D}_i$ periods:
\begin{equation}\label{eq:min-down}
  \sum_{\tau=\max(1,\,t-T^{D}_i+1)}^{t} w_{i,\tau} \;\leq\; 1 - u_{i,t}
  \qquad \forall\, i,\; t
\end{equation}

%-------------------------------------------------------------------
\subsubsection{Constraint 12: Maximum Startup/Shutdown Count (§2.6.3.13)}\label{sec:scuc-max-susd}

\begin{equation}\label{eq:max-starts}
  \sum_{t=1}^{T} v_{i,t} \;\leq\; N^{\max,\mathrm{start}}_i
  \qquad \forall\, i \in \mathcal{G}
\end{equation}
\begin{equation}\label{eq:max-stops}
  \sum_{t=1}^{T} w_{i,t} \;\leq\; N^{\max,\mathrm{stop}}_i
  \qquad \forall\, i \in \mathcal{G}
\end{equation}
where a value of 0 means unlimited. In C++, these constraints are only added
when \texttt{gen\_max\_startups[g] > 0} or \texttt{gen\_max\_shutdowns[g] > 0}.

%-------------------------------------------------------------------
\subsubsection{Constraint 13: Line Flow Limits (§2.6.3.14)}\label{sec:scuc-line}

Using PTDF (Power Transfer Distribution Factor):
\begin{equation}\label{eq:line-flow}
  P^{\min}_l \;\leq\;
  \sum_{i=1}^{N} G_{l,i}\, P_{i,t}
  + \sum_{j=1}^{N_T} G_{l,j}\, T^{O}_{j,t}
  + \sum_{d=1}^{N_{TD}} G_{l,d}\, T^{DC}_{d,t}
  - \sum_{k=1}^{K} G_{l,k}\, D_{k,t}
  + SL^{-}_{l,t}
  \;\leq\;
  P^{\max}_l + SL^{+}_{l,t}
\end{equation}

In the C++ implementation, this is decomposed into two $\leq$ constraints
per line per period, with explicit slack variables:
\begin{align}
  \text{(Fwd):}\quad & \sum_b G_{l,b} \cdot \mathrm{NetInj}_{b,t}
  - SL^{+}_{l,t} \;\leq\; P^{\max}_l \label{eq:line-fwd}\\
  \text{(Rev):}\quad & -\sum_b G_{l,b} \cdot \mathrm{NetInj}_{b,t}
  - SL^{-}_{l,t} \;\leq\; -P^{\min}_l \label{eq:line-rev}
\end{align}
where $\mathrm{NetInj}_{b,t} = \sum_{g \in b} P_{g,t}
      + \sum_{w \in b} P^{W}_{w,t}
      + \sum_{s \in b} P^{PV}_{s,t}
      + \sum_{e \in b}(P^{\mathrm{dis}}_{e,t} - P^{\mathrm{ch}}_{e,t})
      + T^{DC}_{b,t} - D_{b,t}$.

%-------------------------------------------------------------------
\subsubsection{Constraint 14: Section Flow Limits (§2.6.3.15)}\label{sec:scuc-section}

\begin{equation}\label{eq:section-flow}
  P^{\min}_s
  \;\leq\;
  \sum_{i=1}^{N} G_{s,i}\, P_{i,t}
  + \sum_{j=1}^{N_T} G_{s,j}\, T^{O}_{j,t}
  + \sum_{d=1}^{N_{TD}} G_{s,d}\, T^{DC}_{d,t}
  - \sum_{k=1}^{K} G_{s,k}\, D_{k,t}
  + SL^{-}_{s,t}
  \;\leq\;
  P^{\max}_s + SL^{+}_{s,t}
\end{equation}
where $G_{s,k}$ is the aggregate PTDF for section $s$ on bus $k$,
computed as a weighted combination of the PTDFs of lines belonging to the
section. In C++, section PTDFs are precomputed by
\texttt{section\_constraints} and slack variables \texttt{SL\_SEC\_POS/NEG}
are added analogously to line flow.

%-------------------------------------------------------------------
\subsubsection{Constraint 15: Storage (§2.6.3.16)}\label{sec:scuc-storage}

\paragraph{(a) Charge/discharge power limits (§2.6.3.16(1)):}
\begin{align}
  P^{\mathrm{ch,min}}_e \cdot \xi^{+}_{e,t}
    &\;\leq\; P^{\mathrm{ch}}_{e,t} \;\leq\;
    P^{\mathrm{ch,max}}_e \cdot \xi^{+}_{e,t} \label{eq:sto-ch-lim}\\
  P^{\mathrm{dis,min}}_e \cdot \xi^{-}_{e,t}
    &\;\leq\; P^{\mathrm{dis}}_{e,t} \;\leq\;
    P^{\mathrm{dis,max}}_e \cdot \xi^{-}_{e,t} \label{eq:sto-dis-lim}\\
  \xi^{+}_{e,t} + \xi^{-}_{e,t} &\;\leq\; 1 \label{eq:sto-excl}
\end{align}

\paragraph{(b) SOC dynamics (§2.6.3.16(2)):}
\begin{equation}\label{eq:soc-dynamics}
  E_{e,t} = E_{e,t-1}
  + \eta^{\mathrm{ch}}_e \cdot P^{\mathrm{ch}}_{e,t} \cdot \Delta t
  - \frac{P^{\mathrm{dis}}_{e,t}}{\eta^{\mathrm{dis}}_e} \cdot \Delta t
\end{equation}
\begin{equation}\label{eq:soc-limits}
  \underline{E}_{e,t} \;\leq\; E_{e,t} \;\leq\; \bar{E}_{e,t}
  \qquad \forall\, t
\end{equation}

\paragraph{(c) Initial and final SOC (§2.6.3.16(3)):}
\begin{equation}\label{eq:soc-init-final}
  E_{e,0} = E^{\mathrm{ini}}_e, \qquad
  E_{e,T} \geq E^{\mathrm{fin}}_e
\end{equation}

\paragraph{(d) No simultaneous charge/discharge within hour (§2.6.3.16(4)):}
For each hour $n \in [1, 24]$, covering the 4 periods $t \in [4n-3, 4n]$:
\begin{equation}\label{eq:sto-hourly}
  \sum_{t'=4n-3}^{4n} \xi^{+}_{e,t'} \leq 4\,\xi^{+}_{e,n}, \quad
  \sum_{t'=4n-3}^{4n} \xi^{-}_{e,t'} \leq 4\,(1-\xi^{+}_{e,n})
\end{equation}
In the LP relaxation (C++ implementation), this is approximated by:
\begin{equation}\label{eq:sto-simul}
  P^{\mathrm{ch}}_{e,t} + P^{\mathrm{dis}}_{e,t}
  \;\leq\;
  \max\!\bigl(P^{\mathrm{ch,max}}_e,\; P^{\mathrm{dis,max}}_e\bigr)
\end{equation}

\paragraph{(e) Cycle limit (§2.6.3.16(5)):}
\begin{equation}\label{eq:sto-cycle}
  \sum_{t=1}^{T}
  \Bigl(\frac{P^{\mathrm{dis}}_{e,t}}{\eta^{\mathrm{dis}}_e}
  + \eta^{\mathrm{ch}}_e \cdot P^{\mathrm{ch}}_{e,t}\Bigr) \cdot \Delta t
  \;\leq\;
  2 \cdot N^{\mathrm{cycle}}_e \cdot \bar{E}_e
\end{equation}

%-------------------------------------------------------------------
\subsubsection{Constraint 16: DC Line (§2.6.3.17)}\label{sec:scuc-dc}

\paragraph{(a) Multi-terminal DC power balance:}
For a multi-terminal DC group $MT$:
\begin{equation}\label{eq:mt-balance}
  \sum_{s \in J^{\mathrm{send}}_{MT}} T^{DC}_{s,t}
  = \sum_{d \in J^{\mathrm{rec}}_{MT}} T^{DC}_{d,t} \cdot (1 - \rho_d)
\end{equation}
where $\rho_d$ is the loss factor for DC line $d$.

\paragraph{(b) Power limits:}
\begin{equation}\label{eq:dc-limits}
  T^{DC,\min}_{j,t} \;\leq\; T^{DC}_{j,t} \;\leq\; T^{DC,\max}_{j,t}
\end{equation}

\paragraph{(c) Ramp constraints with direction change prohibition:}
\begin{align}
  T^{DC}_{j,t} - T^{DC}_{j,t-1} &\;\leq\; x_{j,t,r} \cdot \Delta P^{+}_j \label{eq:dc-ramp-up}\\
  T^{DC}_{j,t-1} - T^{DC}_{j,t} &\;\leq\; x_{j,t,l} \cdot \Delta P^{-}_j \label{eq:dc-ramp-dn}\\
  x_{j,t,r} + x_{j,t,l} &\;\leq\; 1 \label{eq:dc-ramp-excl}
\end{align}

\paragraph{(d) No reversal between adjacent periods (§2.6.3.17(4)):}
\begin{equation}\label{eq:dc-no-reversal}
  x_{j,t-1,r} + x_{j,t,l} \leq 1, \qquad
  x_{j,t,r} + x_{j,t-1,l} \leq 1
\end{equation}

%-------------------------------------------------------------------
\subsubsection{Constraint 17: Hydro Water Level (§2.6.3.18)}\label{sec:scuc-hydro-wl}

For hydro station $i$ with reservoir:
\begin{equation}\label{eq:hydro-wl}
  Z^{\min}_{i,t}
  \;\leq\;
  Z_{i,0} + \frac{1}{S_i}\sum_{\tau=1}^{t}
  \Bigl(I_{i,\tau} - P_{i,\tau} \cdot h_i - Q^{d}_{i,\tau}
  + P_{\mathrm{up}(i),\tau-s(i)} \cdot h_{\mathrm{up}(i)}
  + Q^{d}_{\mathrm{up}(i),\tau-s(i)}\Bigr)
  \;\leq\;
  Z^{\max}_{i,t}
\end{equation}
where $Z_{i,0}$ is the initial water level, $S_i$ is the reservoir surface
area, $h_i$ is the water consumption rate, $I_{i,\tau}$ is inflow, $Q^{d}_{i}$
is spillage flow, $\mathrm{up}(i)$ is the upstream station, and $s(i)$ is the
upstream delay.

%-------------------------------------------------------------------
\subsubsection{Constraint 18: Hydro Energy Limits (§2.6.3.19)}\label{sec:scuc-hydro-energy}

\begin{equation}\label{eq:hydro-energy}
  Q^{\min}_i \;\leq\;
  \sum_{t=1}^{T} Q_{i,t}
  \;\leq\; Q^{\max}_i
  \qquad \forall\, i \in \mathcal{H}
\end{equation}

%-------------------------------------------------------------------
\subsubsection{Constraint 19: Renewable Output (§2.6.3.20)}\label{sec:scuc-renewable}

For bid-quantity-and-price renewable units:
\begin{equation}\label{eq:re-bounds}
  \alpha_n \cdot P^{\mathrm{st}}_{r,t}
  \;\leq\;
  P_{r,t}
  \;\leq\;
  P^{\mathrm{st}}_{r,t}
  \qquad \forall\, r \in \mathcal{R},\; t \in \mathcal{T}
\end{equation}
$P^{d}_{r,t} = P^{\mathrm{st}}_{r,t} - P_{r,t}$ is the curtailed power.

For bid-quantity-only (price-taker) units:
\begin{equation}\label{eq:re-price-taker}
  P_{r,t} + P^{d}_{r,t} = P^{F}_{r,t}
\end{equation}

%-------------------------------------------------------------------
\subsubsection{Constraint 20: Inter-Provincial Priority Schedule (§2.6.3.21)}\label{sec:scuc-priority}

\begin{equation}\label{eq:priority}
  \sum_{t=1}^{T} Q_{k,t} \;\geq\; Q^{\mathrm{sch}}_k
  \qquad \forall\, k
\end{equation}
where $Q_{k,t}$ is the cleared energy of inter-provincial trading component $k$
and $Q^{\mathrm{sch}}_k$ is the D-2 adjusted priority schedule energy.

%===================================================================
\subsection{SCED -- Security-Constrained Economic Dispatch (§2.6.4)}
\label{sec:sced}
%===================================================================

\subsubsection{SCED Objective (§2.6.4)}

The SCED objective is identical in structure to SCUC but without startup and
minimum-output costs (these are sunk costs given fixed commitment):
\begin{equation}\label{eq:sced-obj}
\begin{aligned}
\min \;\; & \sum_{i=1}^{N}\sum_{t=1}^{T} C_i(P_{i,t})
  \;+\; \sum_{i=1}^{n}\sum_{t=1}^{T} P_{\mathrm{gwf}} \cdot P_{L,i,t} \\
  &+\; M_1 \sum_{l}\sum_{t}(SL^{+}_{l,t}+SL^{-}_{l,t})
  \;+\; M_1 \sum_{s}\sum_{t}(SL^{+}_{s,t}+SL^{-}_{s,t}) \\
  &+\; M_2 \sum_{r}\sum_{t} P^{d}_{r,t}
  \;+\; M_3 \sum_{i}\sum_{t} P^{d,\mathrm{hydro}}_{i,t}
  \;+\; M_4 \sum_{k}\sum_{t} SL_k \\
  &+\; \sum_{e}\sum_{t}
    \bigl(\lambda^{\mathrm{dis}}_e P^{\mathrm{dis}}_{e,t}
    + \lambda^{\mathrm{ch}}_e P^{\mathrm{ch}}_{e,t}\bigr)
\end{aligned}
\end{equation}

\subsubsection{SCED Key Differences from SCUC}

\begin{enumerate}
  \item \textbf{Fixes commitment}: $u_{i,t}$, $v_{i,t}$, $w_{i,t}$ are fixed
        to the SCUC optimal values; the problem becomes a pure LP.
  \item \textbf{Output limits (§2.6.4.5)}: For units that are off in SCUC,
        $P^{\min}_{i,t} = P^{\max}_{i,t} = 0$. For units in startup/shutdown
        transition, limits are fixed to trajectory values.
  \item \textbf{Ramp constraints (§2.6.4.8)}: Simplified---no startup/shutdown
        ramp relaxation terms:
        \begin{equation}\label{eq:sced-ramp}
          |P_{i,t} - P_{i,t-1}| \;\leq\; \Delta P^{U/D}_i
        \end{equation}
  \item All other constraints (balance, reserves, line/section flow, storage,
        DC line, renewables, priority schedule) carry over from SCUC.
\end{enumerate}

\paragraph{C++ implementation.}
In \texttt{solve\_sced()}, the SCUC formulation is rebuilt, then binary
variables (\texttt{IG}, \texttt{SU}, \texttt{SD}) are fixed via their
bounds. The lists \texttt{binary\_idx} and \texttt{integer\_idx} are
cleared so the solver treats the problem as a pure LP.

%===================================================================
\subsection{LMP -- Locational Marginal Price (§2.6.5--2.6.6)}
\label{sec:lmp}
%===================================================================

\subsubsection{LMP Objective (§2.6.5)}

The LMP model has the same objective structure as SCED, with potentially
different penalty factors $M'_1, M'_2, M'_3, M'_4$:
\begin{equation}\label{eq:lmp-obj}
\begin{aligned}
\min \;\; & \sum_{i=1}^{N}\sum_{t=1}^{T} C_i(P_{i,t})
  + \sum_{i=1}^{n}\sum_{t=1}^{T} P_{\mathrm{gwf}} \cdot P_{L,i,t} \\
  &+ M'_1 \sum_l \sum_t (SL^{+}_{l,t}+SL^{-}_{l,t})
  + M'_1 \sum_s \sum_t (SL^{+}_{s,t}+SL^{-}_{s,t}) \\
  &+ M'_2 \sum_r \sum_t P^{d}_{r,t}
  + M'_3 \sum_i \sum_t P^{d,\mathrm{hydro}}_{i,t}
  + M'_4 \sum_k \sum_t SL_k \\
  &+ \sum_e \sum_t (\lambda^{\mathrm{dis}}_e P^{\mathrm{dis}}_{e,t}
    + \lambda^{\mathrm{ch}}_e P^{\mathrm{ch}}_{e,t})
\end{aligned}
\end{equation}

\subsubsection{\texorpdfstring{$\delta$-Neighbourhood Bounds}{delta-Neighbourhood Bounds} (§2.6.5.5)}\label{sec:lmp-delta}

For committed, price-setting generators:
\begin{equation}\label{eq:delta-gen}
  P^{\min,\mathrm{LMP}}_{i,t} = \max\!\bigl((1-\delta)\,P^{\mathrm{SCED}}_{i,t},\;
  P^{\min}_{i,t}\bigr)
\end{equation}
\begin{equation}\label{eq:delta-gen-ub}
  P^{\max,\mathrm{LMP}}_{i,t} = \min\!\bigl((1+\delta)\,P^{\mathrm{SCED}}_{i,t},\;
  P^{\max}_{i,t}\bigr)
\end{equation}

For non-price-setting generators (must-run, in transition, non-market):
\begin{equation}
  P^{\min,\mathrm{LMP}}_{i,t} = P^{\max,\mathrm{LMP}}_{i,t} = P^{\mathrm{SCED}}_{i,t}
\end{equation}

For storage units (§2.6.5.11), analogous $\delta$-bounds apply:
\begin{align}
  P^{\mathrm{ch,max}}_{e,t} &= \min\!\bigl((1+\delta)\,P^{\mathrm{ch,SCED}}_{e,t},\;
    P^{\mathrm{ch,max}}_e\bigr) \\
  P^{\mathrm{dis,max}}_{e,t} &= \min\!\bigl((1+\delta)\,P^{\mathrm{dis,SCED}}_{e,t},\;
    P^{\mathrm{dis,max}}_e\bigr)
\end{align}

\subsubsection{LMP Decomposition (§2.6.6)}\label{sec:lmp-decomp}

The nodal LMP at bus $k$ in period $t$ is:
\begin{equation}\label{eq:lmp-formula}
\boxed{
  \mathrm{LMP}_{k,t} =
  \underbrace{\lambda_t}_{\text{Energy (SMP)}}
  + \underbrace{
    \sum_{l=1}^{N_L}
    \bigl(\mu^{\max}_{l,t} - \mu^{\min}_{l,t}\bigr) \cdot G_{l,k}
    + \sum_{s=1}^{N_S}
    \bigl(\mu^{\max}_{s,t} - \mu^{\min}_{s,t}\bigr) \cdot G_{s,k}
  }_{\text{Congestion component}}
}
\end{equation}
where:
\begin{itemize}[nosep]
  \item $\lambda_t$: Lagrange multiplier of the power balance constraint
        (system marginal price, SMP).
  \item $\mu^{\max}_{l,t}$: multiplier of line $l$ forward flow constraint;
        equals $M'_1$ when the slack variable absorbs flow violation.
  \item $\mu^{\min}_{l,t}$: multiplier of line $l$ reverse flow constraint.
  \item $\mu^{\max}_{s,t}$, $\mu^{\min}_{s,t}$: analogous multipliers for
        section $s$.
  \item $G_{l,k}$: PTDF of bus $k$ on line $l$; $G_{s,k}$: PTDF of bus $k$
        on section $s$.
\end{itemize}

\paragraph{C++ implementation.}
In \texttt{solve\_lmp()}, the LP is solved with a simplex method.
SMP is extracted from the power-balance equality dual.
Congestion components come from line/section flow inequality duals.

%===================================================================
\subsection{Market Clearing Pipeline}\label{sec:pipeline}
%===================================================================

\begin{figure}[H]
\centering
\begin{tikzpicture}[
  box/.style={rectangle, draw=blue!55!black, fill=blue!4, rounded corners=2pt,
              minimum height=0.8cm, text width=4.5cm, align=center},
  arr/.style={-{Stealth[length=3mm]}, thick}
]
  \node[box] (scuc)  {SCUC (MILP)\\[-1pt]{\footnotesize Commitment $u,v,w$}};
  \node[box, below=0.5cm of scuc]  (sced) {SCED (LP)\\[-1pt]{\footnotesize Dispatch $P$}};
  \node[box, below=0.5cm of sced]  (lmp)  {LMP (LP + duals)\\[-1pt]{\footnotesize Nodal prices}};
  \node[box, below=0.5cm of lmp]   (acpf) {AC Power Flow\\[-1pt]{\footnotesize Voltage, losses}};
  \node[box, below=0.5cm of acpf]  (sec)  {Security Check\\[-1pt]{\footnotesize N-1 analysis}};
  \node[box, below=0.5cm of sec]   (stl)  {Settlement\\[-1pt]{\footnotesize Revenue, uplift}};
  \draw[arr] (scuc)  -- node[right,font=\footnotesize]{Fix $u_{i,t}$} (sced);
  \draw[arr] (sced)  -- node[right,font=\footnotesize]{$\delta$-bounds} (lmp);
  \draw[arr] (lmp)   -- (acpf);
  \draw[arr] (acpf)  -- (sec);
  \draw[arr] (sec)   -- (stl);
\end{tikzpicture}
\caption{Day-ahead market clearing pipeline (§2.6.2).}
\label{alg:market-pipeline}
\end{figure}

%===================================================================
\subsection{Summary of C++ Implementation Mapping}
\label{sec:cpp-mapping}
%===================================================================

\begin{longtable}{lll}
\toprule
\textbf{PDF §} & \textbf{Constraint} & \textbf{C++ Location} \\
\midrule
\endfirsthead
\toprule
\textbf{PDF §} & \textbf{Constraint} & \textbf{C++ Location} \\
\midrule
\endhead
2.6.3 Obj    & Objective function           & \texttt{build\_scuc\_formulation}, objective block \\
2.6.3.1      & Power balance                & \texttt{System power balance each period} \\
2.6.3.2      & Positive reserve             & \texttt{Reserve requirements (system + BA)} \\
2.6.3.3      & Negative reserve             & \texttt{Reserve requirements (system + BA)} \\
2.6.3.4      & Primary freq.\ regulation    & \texttt{Spinning reserve constraints} \\
2.6.3.5      & Must-run units               & \texttt{gen\_must\_run} $\Rightarrow$ \texttt{ig.lb = ig.ub = 1} \\
2.6.3.6      & Output expression            & \texttt{PG = Pmin*IG + sum(segments)} \\
2.6.3.7      & Operating cost               & \texttt{lp.c(seg) = price * dt} \\
2.6.3.8      & Output limits                & \texttt{PG <= pmax*IG}, \texttt{PG >= pmin*IG} \\
2.6.3.9      & Group output limits          & \texttt{gen\_group\_constraints} \\
2.6.3.10     & Group energy limits          & \texttt{gen\_group\_energy\_constraints} \\
2.6.3.11     & Ramp rate                    & \texttt{PG - PG\_prev <= ramp\_up + Pmax*SU} \\
2.6.3.12     & Min up/down time             & \texttt{Sum SU $\leq$ IG}, \texttt{Sum SD $\leq$ 1-IG} \\
2.6.3.13     & Max startup/shutdown count   & \texttt{Sum SU $\leq$ gen\_max\_startups} \\
2.6.3.14     & Line flow (PTDF)             & Forward/reverse PTDF + \texttt{SL\_LINE\_POS/NEG} \\
2.6.3.15     & Section flow                 & \texttt{section\_constraints} + \texttt{SL\_SEC\_POS/NEG} \\
2.6.3.16     & Storage                      & SOC dynamics, limits, final SOC, anti-simul, cycle \\
2.6.3.17     & DC line                      & Bounds \& ramp on \texttt{PDC} variables \\
2.6.3.18     & Hydro water level            & \texttt{hydro\_water\_level\_constraints} \\
2.6.3.19     & Hydro energy limits          & \texttt{hydro\_energy\_constraints} \\
2.6.3.20     & Renewable output             & Variable bounds from forecast profiles \\
2.6.3.21     & Priority schedule            & \texttt{priority\_schedule\_constraints} \\
\midrule
2.6.4        & SCED fixes commitment        & \texttt{solve\_sced}: fix IG/SU/SD bounds \\
2.6.5        & LMP $\delta$-neighbourhood   & \texttt{solve\_lmp}: delta bounds on PG/storage \\
2.6.6        & Nodal price formula          & Dual extraction: $\mathrm{SMP} + \sum (\mu^{\max}-\mu^{\min}) \cdot G$ \\
\bottomrule
\end{longtable}

%===================================================================
\subsection{MILP Formulation Tightening: Problem-Specific Cutting Planes}
\label{sec:milp-tightening}
%===================================================================

The SCUC MILP \eqref{eq:scuc-obj}--\eqref{eq:sto-cycle} admits an LP
relaxation whose root gap arises from two distinct structural sources:
\emph{commitment-logic relaxation} (temporal coupling of $u_{i,t}$,
$v_{i,t}$, $w_{i,t}$) and \emph{storage-mode relaxation} (fractional
charging/discharging that violates energy conservation).
This section documents eleven families of valid inequalities that are
separated and appended to the branch-and-cut LP to close these gaps.
All cuts are sparse (number of non-zeros independent of problem size),
certificate-preserving, and compatible with the native solver's Append
Gate and density cap.

%-------------------------------------------------------------------
\subsubsection{LP Relaxation Gap Sources}
\label{sec:lp-gap-sources}
%-------------------------------------------------------------------

\paragraph{Commitment-logic weaknesses.}
\begin{enumerate}[label=(\roman*)]
  \item \textbf{Temporal coupling}: $u_{i,t}$ can take fractional values;
        minimum up/down time constraints are correspondingly weakened.
  \item \textbf{Startup/shutdown logic}: $v_{i,t}$ and $w_{i,t}$ may
        simultaneously be fractional, violating their mutual-exclusion intent.
  \item \textbf{Output--state coupling}: $P_{i,t}$ may be non-zero even when
        $u_{i,t}$ is small, because the product $P^{\max}_i \cdot u_{i,t}$ in
        \eqref{eq:gen-limits} is only linearised at its LP extreme.
  \item \textbf{Startup/shutdown ramp coupling}: the relaxed ramp
        constraints~\eqref{eq:ramp-up}--\eqref{eq:ramp-down} fail to capture
        the multi-period ramp-up trajectory from a cold start.
\end{enumerate}

\paragraph{Storage-mode weaknesses.}
\begin{enumerate}[label=(\roman*)]
  \item \textbf{SOC path infeasibility}: the LP allows $E_{e,t}$ to
        ``jump'' across time periods beyond what is physically achievable.
  \item \textbf{Simultaneous charge/discharge}: $\xi^+_{e,t}=\xi^-_{e,t}=0.5$
        is LP-feasible via \eqref{eq:sto-excl}, creating fictitious energy from
        roundtrip losses ($\eta^{\mathrm{ch}} \eta^{\mathrm{dis}} < 1$).
  \item \textbf{Energy--duration decoupling}: the LP may schedule sustained
        full-power discharge across many periods in excess of usable capacity.
\end{enumerate}

%-------------------------------------------------------------------
\subsubsection{Generator Cut Families (Families 1--6)}
\label{sec:gen-cuts}
%-------------------------------------------------------------------

Let $SU_i$ denote the maximum output achievable in the first period after a
cold start (startup ramp capacity, [MW]), $SD_i$ the maximum output in the
last period before shutdown, and define:
\begin{equation}\label{eq:ramp-horizon}
  K_i \;=\; \left\lceil \frac{P^{\max}_i - SU_i}{\Delta P^U_i} \right\rceil,
  \qquad
  \bar{P}_i^{(k)} \;=\; \min\!\bigl(P^{\max}_i,\; SU_i + k\,\Delta P^U_i\bigr),
  \quad k = 0,1,\ldots,K_i
\end{equation}
$K_i$ is the number of periods required to ramp from cold-start output $SU_i$
to full capacity $P^{\max}_i$.

\paragraph{Family~1: Multi-Period Startup Ramp Inequality.}

For each unit $i \in \mathcal{G}$ and period $t \in \mathcal{T}$:
\begin{equation}\label{eq:cut-ramp-up}
  \boxed{
    P_{i,t}
    \;\leq\;
    P^{\max}_i \cdot u_{i,t}
    \;-\;
    \sum_{k=1}^{\min(K_i,\,t)} \bigl(P^{\max}_i - \bar{P}_i^{(k)}\bigr)\,v_{i,t-k+1}
  }
\end{equation}
\textit{Validity}: if $u_{i,t}=1$ and unit started $k-1$ periods ago
($v_{i,t-k+1}=1$), the right-hand side equals $\bar{P}_i^{(k)}$, which is
precisely the physical ramp-up ceiling after $k$ periods.

\paragraph{Family~2: Startup Clique and Strengthened Interval Up Inequality.}

\emph{Startup clique} (no two startups within one up-time window):
\begin{equation}\label{eq:cut-startup-clique}
  \sum_{t'=t}^{t+T^U_i-1} v_{i,t'} \;\leq\; 1
  \qquad \forall\, i,\; t
\end{equation}

\emph{Strengthened interval up inequality}: for any sub-window of length
$L \leq T^U_i$ starting at $t_1$:
\begin{equation}\label{eq:cut-interval-up}
  \boxed{
    \sum_{t'=t_1}^{t_1+L-1} u_{i,t'}
    \;\geq\;
    \sum_{k=0}^{L-1} (L - k)\,v_{i,t_1+k}
  }
\end{equation}
\textit{Validity}: if a startup at $t_1+k$ occurs (integer solution), then
$u_{i,t'}=1$ for all $t' \in [t_1+k,\,t_1+L-1]$, contributing exactly
$L-k$ to the left-hand sum.  The symmetric shutdown clique and interval-down
inequality hold by replacing $v_{i,t'}$ with $w_{i,t'}$ and $u_{i,t'}$ with
$1-u_{i,t'}$.

\paragraph{Family~3: Multi-Period Shutdown Ramp-Down Inequality.}

Let $\bar{P}_i^{\mathrm{dn},(k)} = \min(P^{\max}_i,\, SD_i + k\,\Delta P^D_i)$ and
$K_i^{\mathrm{dn}} = \lceil(P^{\max}_i - SD_i)/\Delta P^D_i\rceil$.
\begin{equation}\label{eq:cut-ramp-down}
  \boxed{
    P_{i,t}
    \;\leq\;
    P^{\max}_i \cdot u_{i,t}
    \;-\;
    \sum_{k=1}^{\min(K_i^{\mathrm{dn}},\, T-t)} \bigl(P^{\max}_i - \bar{P}_i^{\mathrm{dn},(k)}\bigr)\,w_{i,t+k}
  }
\end{equation}
This is the time-reverse of Family~1: if a shutdown will occur $k$ periods
ahead, the current output is bounded by the pre-shutdown ramp-down ceiling.

\paragraph{Family~4: Multi-Level Startup Cost Tightening.}

For units with state-dependent startup costs (hot/warm/cold), let
$\tau_i^{\mathrm{hot}}$ and $\tau_i^{\mathrm{warm}}$ be the hot- and warm-start
cooling thresholds (periods), $DT_i \leq \tau_i^{\mathrm{hot}} \leq \tau_i^{\mathrm{warm}}$.
Introduce binary startup-type indicators $\delta_{i,t}^{\mathrm{warm}}$ and
$\delta_{i,t}^{\mathrm{cold}}$:
\begin{align}
  v_{i,t} &\;\leq\; \sum_{t'=t-\tau_i^{\mathrm{hot}}}^{t-T^D_i} u_{i,t'}
                   + \delta_{i,t}^{\mathrm{warm}} + \delta_{i,t}^{\mathrm{cold}}
  \label{eq:cut-sutype-a}\\
  \delta_{i,t}^{\mathrm{warm}} + \delta_{i,t}^{\mathrm{cold}}
  &\;\leq\; 1 - \sum_{t'=t-\tau_i^{\mathrm{hot}}}^{t-T^D_i} u_{i,t'}
  \label{eq:cut-sutype-b}\\
  \delta_{i,t}^{\mathrm{cold}}
  &\;\leq\; 1 - \sum_{t'=t-\tau_i^{\mathrm{warm}}}^{t-\tau_i^{\mathrm{hot}}-1} u_{i,t'}
  \label{eq:cut-sutype-c}
\end{align}
These inequalities force the LP to identify the startup type correctly,
preventing fractional mixing of hot and cold costs.

\paragraph{Family~5: Three-Period Window Inequality.}

Simultaneously tightening the ramp-up, ramp-down, and commitment logic:
\begin{equation}\label{eq:cut-3period}
  \boxed{
    P_{i,t}
    \;\leq\;
    P^{\max}_i \cdot u_{i,t}
    \;-\; (P^{\max}_i - SU_i)\,v_{i,t}
    \;-\; (P^{\max}_i - SD_i)\,w_{i,t+1}
  }
\end{equation}
\textit{Interpretation}: if the unit starts in period $t$ it cannot exceed $SU_i$;
if it shuts down in period $t+1$ it must be no greater than $SD_i$.

\paragraph{Family~6: Symmetry-Breaking Inequalities.}

For a group $\mathcal{G}^{\mathrm{sym}} = \{i_1, i_2, \ldots, i_K\}$ of
identical units (same $P^{\max}$, $P^{\min}$, $T^U$, $T^D$, cost curves):

\emph{Lexicographic ordering}: commit units in label order,
\begin{equation}\label{eq:cut-sym-lex}
  u_{i_j,t} \;\geq\; u_{i_{j+1},t}
  \qquad \forall\, j=1,\ldots,K-1,\quad \forall\, t
\end{equation}

\emph{Cumulative startup ordering} (stronger):
\begin{equation}\label{eq:cut-sym-cumstart}
  \sum_{t'=1}^{t} v_{i_j,t'} \;\geq\; \sum_{t'=1}^{t} v_{i_{j+1},t'}
  \qquad \forall\, j,\; t
\end{equation}
Symmetry-breaking cuts are generated once at the root node and remain
globally valid throughout the branch-and-cut tree.

%-------------------------------------------------------------------
\subsubsection{Storage Cut Families (Families 7--11)}
\label{sec:storage-cuts}
%-------------------------------------------------------------------

The original six generator cuts remain valid in the presence of storage
(they only involve $\mathcal{G}$ variables), but they cannot close the
additional LP gap introduced by $\mathcal{E}$ variables.
The following five families address the storage-specific sources identified
in Section~\ref{sec:lp-gap-sources}.

\paragraph{Family~7: SOC Reachability Inequality.}

For each storage unit $e \in \mathcal{E}$ and any interval $[t_1, t_2]$:
\begin{align}
  E_{e,t_2} - E_{e,t_1}
  &\;\leq\;
  \eta^{\mathrm{ch}}_e \cdot P^{\mathrm{ch,max}}_e \cdot \Delta t
  \sum_{t=t_1+1}^{t_2} \xi^{+}_{e,t}
  \label{eq:cut-soc-up}\\
  E_{e,t_1} - E_{e,t_2}
  &\;\leq\;
  \frac{P^{\mathrm{dis,max}}_e}{\eta^{\mathrm{dis}}_e} \cdot \Delta t
  \sum_{t=t_1+1}^{t_2} \xi^{-}_{e,t}
  \label{eq:cut-soc-dn}
\end{align}
\textit{Validity} (charging direction, \eqref{eq:cut-soc-up}):
from the SOC dynamics \eqref{eq:soc-dynamics},
$E_{e,t_2}-E_{e,t_1} \leq \sum_{t} \eta^{\mathrm{ch}}_e P^{\mathrm{ch}}_{e,t} \cdot\Delta t
\leq \sum_{t} \eta^{\mathrm{ch}}_e P^{\mathrm{ch,max}}_e \cdot \xi^{+}_{e,t} \cdot \Delta t$.
The LP may violate \eqref{eq:cut-soc-up} when $\xi^+$ is fractional and
the product $P^{\mathrm{ch,max}}_e \cdot \xi^+_{e,t}$ over-estimates achievable charging.
The window length is capped at $W \leq 12$ periods in the separator to
control the number of cut candidates.

\paragraph{Family~8: Simultaneous Charge/Discharge Disjunctive Cut.}

Beyond the period-level exclusion \eqref{eq:sto-excl}, the following cut
exploits the energy loss implied by simultaneous operation:
\begin{equation}\label{eq:cut-cd-disj}
  \eta^{\mathrm{ch}}_e \cdot P^{\mathrm{ch}}_{e,t} \cdot \Delta t
  +
  \frac{P^{\mathrm{dis}}_{e,t}}{\eta^{\mathrm{dis}}_e} \cdot \Delta t
  \;\leq\;
  \max\!\Bigl(\eta^{\mathrm{ch}}_e P^{\mathrm{ch,max}}_e,\;
              \tfrac{P^{\mathrm{dis,max}}_e}{\eta^{\mathrm{dis}}_e}\Bigr)
  \cdot (\xi^{+}_{e,t} + \xi^{-}_{e,t}) \cdot \Delta t
\end{equation}
In an integer solution $\xi^+_{e,t}+\xi^-_{e,t} \in \{0,1\}$, so the right-hand
side equals either 0 (idle) or the larger of the two physical power rates.
When the LP places both indicators at $0.5$, the right-hand side becomes
$\max(\cdot)$, which is tighter than the sum of the individual big-$M$
bounds given by \eqref{eq:sto-ch-lim}--\eqref{eq:sto-dis-lim}.

\paragraph{Family~9: Energy--Duration Coupling Inequality.}

For any discharge window $[t_1, t_2]$:
\begin{equation}\label{eq:cut-energy-duration}
  \boxed{
    \sum_{t=t_1}^{t_2} \frac{P^{\mathrm{dis}}_{e,t}}{\eta^{\mathrm{dis}}_e} \cdot \Delta t
    \;\leq\;
    (E_{e,t_1} - \underline{E}_e)
    \;+\;
    \eta^{\mathrm{ch}}_e P^{\mathrm{ch,max}}_e \cdot \Delta t
    \sum_{t=t_1+1}^{t_2} \xi^{+}_{e,t}
  }
\end{equation}
\textit{Interpretation}: the total energy extracted over the window cannot
exceed the available energy at $t_1$ above the minimum SOC, augmented by
any charging that occurs during the same window.

\paragraph{Family~10: Cross-Device Substitution Inequality (Cover Type).}

For each period $t$ and any minimal cover set
$\mathcal{C}_t \subseteq \mathcal{G}^{a} \cup \mathcal{E}^{a}$ satisfying

$$
  \sum_{i \in \mathcal{C}_t \cap \mathcal{G}} P^{\max}_i
  + \sum_{e \in \mathcal{C}_t \cap \mathcal{E}} P^{\mathrm{dis,max}}_e
  > D^a_t - \sum_{e \notin \mathcal{C}_t} P^{\mathrm{dis,max}}_e,
$$

the resulting cover inequality is:
\begin{equation}\label{eq:cut-cover}
  \sum_{i \in \mathcal{C}_t \cap \mathcal{G}} u_{i,t}
  + \sum_{e \in \mathcal{C}_t \cap \mathcal{E}} \xi^{-}_{e,t}
  \;\leq\; |\mathcal{C}_t| - 1
\end{equation}
This forbids every device in the cover from simultaneously operating at or
near full output, which is unnecessary to meet demand.  In practice, the
separator identifies the most-violated cover of size $|\mathcal{C}_t| \leq 5$
at the root node only, limiting the $\mathcal{O}(|\mathcal{G}||\mathcal{E}|)$
enumeration to tractable runs.

\paragraph{Family~11: Cyclic SOC Feasibility Inequality.}

When the problem requires $E_{e,T} \geq E^{\mathrm{fin}}_e \approx E^{\mathrm{ini}}_e$
(daily cycling), summing \eqref{eq:soc-dynamics} over all periods gives:
\begin{equation}\label{eq:soc-balance}
  \sum_{t=1}^{T} \eta^{\mathrm{ch}}_e P^{\mathrm{ch}}_{e,t} \cdot \Delta t
  \;\geq\;
  \sum_{t=1}^{T} \frac{P^{\mathrm{dis}}_{e,t}}{\eta^{\mathrm{dis}}_e} \cdot \Delta t
  + (E^{\mathrm{fin}}_e - E^{\mathrm{ini}}_e)
\end{equation}
The tightened form replaces $P^{\mathrm{ch}}_{e,t}$ by its binary-state upper
bound, yielding a cut on the charging indicator sum:
\begin{equation}\label{eq:cut-cyclic-soc}
  \boxed{
    \eta^{\mathrm{ch}}_e P^{\mathrm{ch,max}}_e \cdot \Delta t
    \sum_{t=1}^{T} \xi^{+}_{e,t}
    \;\geq\;
    \sum_{t=1}^{T} \frac{P^{\mathrm{dis}}_{e,t}}{\eta^{\mathrm{dis}}_e} \cdot \Delta t
    + (E^{\mathrm{fin}}_e - E^{\mathrm{ini}}_e)
  }
\end{equation}
This global cut is generated once per storage unit at the root node.

%-------------------------------------------------------------------
\subsubsection{Cut Properties and Separation Complexity}
\label{sec:cut-properties}
%-------------------------------------------------------------------

\begin{longtable}{llllll}
\toprule
\textbf{Family} & \textbf{Description} & \textbf{Max candidates} &
  \textbf{Non-zeros} & \textbf{Separation cost} & \textbf{Depth limit} \\
\midrule
\endfirsthead
\toprule
\textbf{Family} & \textbf{Description} & \textbf{Max candidates} &
  \textbf{Non-zeros} & \textbf{Separation cost} & \textbf{Depth limit} \\
\midrule
\endhead
1 & Startup ramp   & $O(NT)$        & $\leq K_i{+}2$        & $O(NTK_{\max})$     & 0--10 \\
2 & Interval up    & $O(NT\cdot T^U_{\max})$ & $\leq 2T^U_{\max}$ & $O(NT(T^U_{\max})^2)$ & 0--10 \\
3 & Shutdown ramp  & $O(NT)$        & $\leq K_i^{\mathrm{dn}}{+}2$ & $O(NTK^{\mathrm{dn}}_{\max})$ & 0--10 \\
4 & Startup cost   & $O(NT)$        & $\leq 2\tau_{\max}$   & $O(NT\tau_{\max})$  & 0 only \\
5 & 3-period window & $O(NT)$       & 3                     & $O(NT)$             & 0--10 \\
6 & Symmetry break & $O(N_{\mathrm{sym}}T)$ & 2          & $O(N_{\mathrm{sym}}T)$ & 0 only \\
\midrule
7 & SOC reachability & $O(N_{ES}TW)$ & $\leq W{+}2$         & $O(N_{ES}TW^2)$     & 0--10 \\
8 & Charge/discharge & $O(N_{ES}T)$  & 4                    & $O(N_{ES}T)$        & 0--10 \\
9 & Energy--duration & $O(N_{ES}TW)$ & $\leq W{+}2$         & $O(N_{ES}TW)$       & 0--10 \\
10 & Cross-device   & $O(T\cdot 2^{|\mathcal{C}|})$ & $|\mathcal{C}|$  & $O(T\cdot NG\cdot N_{ES})$ & 0 only \\
11 & Cyclic SOC     & $O(N_{ES})$   & $\leq T{+}1$          & $O(N_{ES}T)$        & 0 only \\
\bottomrule
\end{longtable}

$W \leq 12$ is the window-length cap for SOC cuts.  All families satisfy the
native solver's density cap ($\leq 0.05 \cdot n$ non-zeros) because non-zero
counts grow with machine parameters ($K_i$, $T^U_i$, $W$), not with total
problem size~$n$.

%-------------------------------------------------------------------
\subsubsection{Layered Separation Strategy}
\label{sec:cut-strategy}
%-------------------------------------------------------------------

Cuts are generated by a \texttt{UCSeparationOracle} (see
\texttt{src/engine/branch\_and\_cut.cpp}) that dispatches to per-family
separators according to tree depth:

\begin{enumerate}
  \item \textbf{Root node ($d=0$)}: Families~4, 6, 10, 11 are generated
        as pre-formulation rows and never revisited.
  \item \textbf{Shallow nodes ($1 \leq d \leq 10$)}: Families~1, 2, 3, 5,
        7, 8, 9 are active as B\&B separation cuts.
  \item \textbf{Deep nodes ($d > 10$)}: no problem-specific cuts; the
        solver relies on domain propagation and reduced-cost fixing.
\end{enumerate}

\paragraph{Pre-formulation vs.~B\&B separation.}
Families~6 and 11 are the only generic families appended to the
LP \emph{before} branch-and-cut begins.  Families~1, 2, 3, and~5 mix
continuous dispatch outputs with startup/shutdown binaries
($v_{i,t'}$, $w_{i,t'}$).  In the LP relaxation these binaries are
fractional, which can drive the right-hand side of
\eqref{eq:cut-ramp-up}--\eqref{eq:cut-3period} below zero
while the lower bound $P^{\min}_i u_{i,t} > 0$ is simultaneously active,
producing an infeasible LP.  They must therefore be added lazily
at integer nodes where $v_{i,t'} \in \{0,1\}$ is already fixed.

After each round of cut generation, violated cuts are filtered by an
efficacy threshold $\epsilon_{\mathrm{cut}} \geq 10^{-4}$ and a maximum
density cap before being appended via \texttt{AppendCuts()}.
Numerical safety for Families~1 and 3 (coefficients proportional to
$P^{\max}_i$) is enforced by normalising each cut row by $P^{\max}_i$
before appending, so all coefficients lie in $[0,1]$.

\paragraph{Interaction with RENS heuristic.}
Problem-specific cuts tighten the LP solution, increasing the fraction of
$u_{i,t}$, $v_{i,t}$, $w_{i,t}$, and $\xi^{\pm}_{e,t}$ variables that are
already integer-valued in the LP optimum.  The Residual RENS heuristic
(§\ref{sec:native-branch-and-cut}) then fixes this larger integer-valued
set, reduces the residual sub-problem size, and finds an incumbent
substantially faster.  The expected effect is:
\begin{equation}
  \underbrace{\text{Tighter LP}}_{\text{Cuts}}
  \;\xrightarrow{\text{more fixed }u/\xi}\;
  \underbrace{\text{Smaller RENS sub-problem}}_{\text{faster incumbent}}
  \;\xrightarrow{\text{small gap}}\;
  \underbrace{\text{Fewer proof nodes}}_{\text{faster termination}}
\end{equation}

%===================================================================
\subsection{Market-Model-Specific Cutting Planes (SCUC Extensions)}
\label{sec:scuc-market-cuts}
%===================================================================

The eleven generic cut families in Section~\ref{sec:milp-tightening} apply
to any UC formulation.  The Southern-China SCUC model
(Section~\ref{sec:scuc}) exposes additional combinatorial structure that
enables nine further families (A--I).  These families exploit features
unique to this market model: piecewise-linear bid curves, multi-type
reserve requirements, monitored sections, intra- and inter-regional DC
interconnections, hydro reservoirs, inter-provincial priority schedules,
and a 96-period (15\,min) time resolution.

%-------------------------------------------------------------------
\subsubsection{Family A: Bid-Segment Cumulative Coupling}
\label{sec:cut-A}
%-------------------------------------------------------------------

The output decomposition \eqref{eq:gen-output} and segment bounds
\eqref{eq:seg-bounds} only enforce each segment $m$ individually.
Because bid prices are non-decreasing ($C_{i,1} \leq C_{i,2} \leq
\cdots \leq C_{i,N_M}$), dispatch fills segments sequentially: segment
$m$ is used only after segment $m-1$ is saturated.  The LP relaxation
violates this ordering when $u_{i,t}$ is fractional.

\paragraph{Cumulative segment cut.}
For each unit $i$, period $t$, and prefix $k = 1, \ldots, N_M$:
\begin{equation}\label{eq:cut-A-cumul}
  \boxed{
    \sum_{m=1}^{k} \Delta P_{i,t,m}
    \;\leq\;
    \Bigl(\sum_{m=1}^{k} \overline{\Delta P}_{i,m}\Bigr) \cdot u_{i,t}
  }
\end{equation}
This is tighter than the individual bounds $\Delta P_{i,t,m} \leq
\overline{\Delta P}_{i,m}\,u_{i,t}$ because it requires the \emph{total}
committed capacity of the first $k$ segments to respect the binary state,
preventing the LP from spreading fractional commitment uniformly across
all segments.

%-------------------------------------------------------------------
\subsubsection{Family B: Reserve--Commitment Cover Inequality}
\label{sec:cut-B}
%-------------------------------------------------------------------

The positive-reserve constraint \eqref{eq:pos-reserve} and the
reserve-coupling bound \eqref{eq:reserve-coupling} jointly require
sufficient committed capacity in each balancing area $a$.  The LP
relaxation allows fractional $u_{i,t}$ to provide the required head-room
at a fraction of the integer cost.

\paragraph{Minimum-commitment cover for reserves.}
For each area $a$ and period $t$, let $\mathcal{F}^a_t$ be the set of
units that cannot individually cover the residual reserve requirement
after subtracting storage and import.  A minimum cover
$\mathcal{C}^a_t \subseteq \mathcal{G}^a$ satisfies

\begin{equation}\label{eq:cut-B-cover-def}
  \sum_{i \in \mathcal{G}^a \setminus \mathcal{C}^a_t}
    (P^{\max}_i - P^{\min}_i)
  \;<\; R^{U,a}_t + D^a_t
    - \sum_{e \in \mathcal{E}^a} P^{\mathrm{dis,max}}_e
    - T^{\mathrm{import,max},a}
\end{equation}

yielding the reserve cover inequality:
\begin{equation}\label{eq:cut-B-cover}
  \boxed{
    \sum_{i \in \mathcal{G}^a \setminus \mathcal{C}^a_t} u_{i,t}
    \;\geq\;
    \biggl\lceil
      \frac{R^{U,a}_t - \sum_{e}P^{\mathrm{dis,max}}_e
            - \sum_{i \in \mathcal{C}^a_t}(P^{\max}_i-P^{\min}_i)}
           {\max_{i \notin \mathcal{C}^a_t}(P^{\max}_i - P^{\min}_i)}
    \biggr\rceil
  }
\end{equation}

\paragraph{Reserve--ramp startup tightening.}
A unit that just started ($v_{i,t}=1$) has a reduced upward ramp
capability.  The achievable spinning reserve at startup is bounded by:
\begin{equation}\label{eq:cut-B-ramp}
  \boxed{
    R^{\mathrm{spin}}_{i,t}
    \;\leq\;
    \min(\Delta P^U_i,\; P^{\max}_i - P^{\min}_i)\,u_{i,t}
    \;-\; (P^{\max}_i - P^{\min}_i - \Delta P^U_i)^{\!+} v_{i,t}
  }
\end{equation}

%-------------------------------------------------------------------
\subsubsection{Family C: Monitored-Section Commitment Inequality}
\label{sec:cut-C}
%-------------------------------------------------------------------

Section flow constraints \eqref{eq:section-flow} aggregate multiple
transmission lines.  Cutting at the section level produces fewer, denser
cuts that are more effective than per-line cuts for large networks.

Define the generation shift factor for section $s$ and unit $i$ as:
\begin{equation}
  \mathrm{GSF}_{s,i} = \sum_{l \in \mathrm{lines}(s)} G_{l,\,\mathrm{bus}(i)}
\end{equation}

\paragraph{Forward section commitment cut.}
For the forward limit of section $s$ in period $t$:
\begin{equation}\label{eq:cut-C-fwd}
  \boxed{
    \sum_{\substack{i \in \mathcal{G}: \\ \mathrm{GSF}_{s,i} > 0}}
      \mathrm{GSF}_{s,i} \cdot P^{\max}_i \cdot u_{i,t}
    \;\geq\;
    P^{\max}_s + \mathrm{NetDemand}^{s,+}_t
  }
\end{equation}
where $\mathrm{NetDemand}^{s,+}_t$ is the section-facing net demand
(load minus storage dispatch and imports on the forward side).
An analogous reverse-direction cut covers the lower section limit.

%-------------------------------------------------------------------
\subsubsection{Family D: DC-Line Direction--Commitment Coupling}
\label{sec:cut-D}
%-------------------------------------------------------------------

The DC line model (Section~\ref{sec:scuc-dc}) prohibits direction reversal
between adjacent periods and limits the ramp.  Introducing a binary
direction indicator $\delta_{d,t} \in \{0,1\}$ (1 = forward) gives:
\begin{align}
  T^{DC}_{d,t} &\leq T^{DC,\max}_d \cdot \delta_{d,t} \label{eq:dc-dir-ub}\\
  T^{DC}_{d,t} &\geq T^{DC,\min}_d \cdot (1-\delta_{d,t}) \label{eq:dc-dir-lb}
\end{align}

\paragraph{Receiving-end commitment requirement.}
When the DC link does not deliver to area $a_{\mathrm{recv}}$
($\delta_{d,t} = 0$), local generation must cover the full load plus reserve:
\begin{equation}\label{eq:cut-D-recv}
  \boxed{
    \sum_{i \in \mathcal{G}^{a_{\mathrm{recv}}}} P^{\max}_i \cdot u_{i,t}
    \;\geq\;
    D^{a_{\mathrm{recv}}}_t + R^{U,a_{\mathrm{recv}}}_t
    - T^{DC,\max}_d \cdot \delta_{d,t}
  }
\end{equation}

\paragraph{DC direction persistence cut (analogue of min-up-time).}
Once a DC direction is established, forcing it to persist for $L$ periods
reduces LP oscillation:
\begin{equation}\label{eq:cut-D-persist}
  \sum_{t'=t}^{t+L-1} \delta_{d,t'} \;\geq\; L \cdot (\delta_{d,t} - \delta_{d,t-1})^+
  \qquad L = 2, 3, 4
\end{equation}

%-------------------------------------------------------------------
\subsubsection{Family E: Storage Cycle-Limit Strengthening}
\label{sec:cut-E}
%-------------------------------------------------------------------

The cycle-energy bound \eqref{eq:sto-cycle} limits total throughput but
does not constrain the number of discharge periods, allowing the LP to
exploit fractional $\xi^-_{e,t}$.

\paragraph{Discharge-period count bound.}
Each full charge--discharge cycle occupies at least $2$ active periods.
Hence the total number of discharge periods is bounded:
\begin{equation}\label{eq:cut-E-count}
  \boxed{
    \sum_{t=1}^{T} \xi^-_{e,t} \;\leq\; 2\,N^{\mathrm{cycle}}_e
  }
\end{equation}

\paragraph{Window cycle quota cut.}
For any sub-window $[t_1, t_2]$ of length $L$:
\begin{equation}\label{eq:cut-E-window}
  \boxed{
    \sum_{t=t_1}^{t_2} \frac{P^{\mathrm{dis}}_{e,t}}{\eta^{\mathrm{dis}}_e} \cdot \Delta t
    \;\leq\;
    N^{\mathrm{cycle}}_e \cdot \frac{L}{T} \cdot (\bar{E}_e - \underline{E}_e)
    + (E_{e,t_1} - \underline{E}_e)
  }
\end{equation}
The right-hand side is the pro-rated cycle quota for the window plus the
energy already available at the start of the window.

%-------------------------------------------------------------------
\subsubsection{Family F: Hydro--Thermal Substitution Inequality}
\label{sec:cut-F}
%-------------------------------------------------------------------

Because hydro generation can substitute for thermal commitment, the LP
over-exploits fractional hydro flexibility to defer binary startups.
Let $E^{\mathrm{total}}_h$ be the total available hydro energy of
station $h$ over the horizon, and $\mathcal{T}_{\mathrm{peak}}$ the set
of peak-load periods.

\paragraph{Thermal capacity sufficiency cut.}
For each peak period $t \in \mathcal{T}_{\mathrm{peak}}$:
\begin{equation}\label{eq:cut-F-cap}
  \boxed{
    \sum_{i \in \mathcal{G}^a} P^{\max}_i \cdot u_{i,t}
    \;\geq\;
    D^a_t + R^{U,a}_t
    - \sum_{h \in \mathcal{H}^a} P^{\max}_h
    - \sum_{e \in \mathcal{E}^a} P^{\mathrm{dis,max}}_e
    - T^{\mathrm{import,max},a}
  }
\end{equation}

\paragraph{Hydro-energy--thermal minimum-output coupling.}
The sum of thermal minimum output over peak periods cannot exceed the
demand surplus after hydro and imports are accounted for:
\begin{equation}\label{eq:cut-F-energy}
  \sum_{t \in \mathcal{T}_{\mathrm{peak}}}
    \sum_{i \in \mathcal{G}^a} P^{\min}_i \cdot u_{i,t}
  \;\leq\;
  \sum_{t \in \mathcal{T}_{\mathrm{peak}}} D^a_t
  - \sum_{h \in \mathcal{H}^a} E^{\mathrm{total}}_h
  - \sum_t T^{\mathrm{import},a}_t
\end{equation}

%-------------------------------------------------------------------
\subsubsection{Family G: Extended Time-Window Clique (96-Period)}
\label{sec:cut-G}
%-------------------------------------------------------------------

With $T = 96$ periods and minimum up- plus down-times $T^U_i + T^D_i$ that
may span 16--48 periods, the extended window clique is substantially tighter
than the standard startup clique \eqref{eq:cut-startup-clique}.

\paragraph{$(T^U_i + T^D_i)$-period startup clique.}
No two startups can occur within a window that spans one full on--off cycle:
\begin{equation}\label{eq:cut-G-clique}
  \boxed{
    \sum_{t'=t}^{t+T^U_i+T^D_i-1} v_{i,t'} \;\leq\; 1
    \qquad \forall\, i,\; t = 1,\ldots,T-(T^U_i+T^D_i)+1
  }
\end{equation}

\paragraph{Total startup count bound.}
Over the full horizon, the maximum number of startups is:
\begin{equation}\label{eq:cut-G-maxstart}
  \sum_{t=1}^{T} v_{i,t}
  \;\leq\;
  \left\lfloor \frac{T}{T^U_i + T^D_i} \right\rfloor
  \qquad \forall\, i
\end{equation}

\paragraph{Strengthened 96-period interval-up inequality.}
For windows of length $W = T^U_i + T^D_i$, the interval-up inequality
\eqref{eq:cut-interval-up} strengthens to:
\begin{equation}\label{eq:cut-G-interval}
  \sum_{t'=t}^{t+W-1} u_{i,t'}
  \;\geq\;
  T^U_i \cdot v_{i,t}
  + T^U_i \cdot v_{i,t+T^U_i+T^D_i} \cdot \mathbf{1}[t+T^U_i+T^D_i \leq T]
\end{equation}
requiring at least $T^U_i$ consecutive on-periods for every startup within
the window.

%-------------------------------------------------------------------
\subsubsection{Family H: Inter-Provincial Schedule--Commitment Coupling}
\label{sec:cut-H}
%-------------------------------------------------------------------

The priority schedule constraint \eqref{eq:priority} is enforced with a
high-penalty slack $SL_k$.  In integer solutions, meeting a large export
schedule may require additional committed capacity beyond what the LP
fractional solution uses.

\paragraph{Export-commitment capacity cut.}
For each area $a$ and period $t$ with net export plan $T^{\mathrm{plan,exp},a}_t$:
\begin{equation}\label{eq:cut-H-cap}
  \boxed{
    \sum_{i \in \mathcal{G}^a} P^{\max}_i \cdot u_{i,t}
    + \sum_{e \in \mathcal{E}^a} P^{\mathrm{dis,max}}_e \cdot \xi^-_{e,t}
    \;\geq\;
    D^a_t + T^{\mathrm{plan,exp},a}_t + R^{U,a}_t
  }
\end{equation}

\paragraph{Export-ramp startup cut.}
When the planned export ramps up by $\Delta T^{\mathrm{plan},a}_t$ between
periods $t-1$ and $t$:
\begin{equation}\label{eq:cut-H-ramp}
  \sum_{i \in \mathcal{G}^a} \Delta P^U_i \cdot u_{i,t-1}
  + \sum_{i \in \mathcal{G}^a} SU_i \cdot v_{i,t}
  \;\geq\; \Delta T^{\mathrm{plan},a}_t
\end{equation}

%-------------------------------------------------------------------
\subsubsection{Family I: Primary Frequency Regulation--Storage Joint Cut}
\label{sec:cut-I}
%-------------------------------------------------------------------

The PFR requirement \eqref{eq:pf-reserve} must be met by thermal and hydro
units.  When storage provides PFR, its contribution is limited by the
available SOC headroom above $\underline{E}_e$.

\paragraph{PFR--SOC joint cut.}
Let $\Delta t_{\mathrm{pf}}$ be the PFR response window [hours] and
$P^{\mathrm{pf,max}}_e = \min(P^{\mathrm{dis,max}}_e,\; (\bar{E}_e -
\underline{E}_e)/(2\,\Delta t_{\mathrm{pf}}))$:
\begin{equation}\label{eq:cut-I-pfr}
  \boxed{
    \sum_{i \in \mathcal{G}} R^{\mathrm{pf,max}}_i \cdot u_{i,t}
    + \sum_{e \in \mathcal{E}} P^{\mathrm{pf,max}}_e \cdot \xi^-_{e,t}
    \;\geq\; R^{\mathrm{pf}}_t
    \qquad \forall\, t
  }
\end{equation}

\paragraph{SOC-dependent PFR cut.}
When the LP SOC variable is available:
\begin{equation}\label{eq:cut-I-soc}
  \sum_{i \in \mathcal{G}} R^{\mathrm{pf,max}}_i \cdot u_{i,t}
  + \sum_{e \in \mathcal{E}}
      \frac{E_{e,t} - \underline{E}_e}{\Delta t_{\mathrm{pf}}}
  \;\geq\; R^{\mathrm{pf}}_t
\end{equation}
This is a convex valid inequality: it tightens the PFR feasibility set
by using the actual SOC value rather than the design maximum.

%-------------------------------------------------------------------
\subsubsection{Market-Cut Priority and Separation Strategy}
\label{sec:market-cut-strategy}
%-------------------------------------------------------------------

\begin{longtable}{lllllll}
\toprule
\textbf{Family} & \textbf{Description} & \textbf{Candidates} &
  \textbf{Non-zeros} & \textbf{Expected gap closing} & \textbf{Depth limit} & \textbf{Impl.} \\
\midrule
\endfirsthead
\toprule
\textbf{Family} & \textbf{Description} & \textbf{Candidates} &
  \textbf{Non-zeros} & \textbf{Expected gap closing} & \textbf{Depth limit} & \textbf{Impl.} \\
\midrule
\endhead
G & Extended clique      & $O(NT)$             & $\leq T^U_i{+}T^D_i$ & 25--30\% & 0--10 & \checkmark \\
B & Reserve--commitment  & $O(|\mathcal{A}|T)$ & $\leq N^a$            & 20--25\% & 0--5  & \checkmark \\
C & Section commitment   & $O(N_S T)$          & $\leq N$              & 15--20\% & 0--5  & \checkmark \\
A & Bid-segment cumul.   & $O(N N_M T)$        & $\leq N_M{+}1$        & 10--15\% & 0--5  & \checkmark \\
E & Cycle strengthening  & $O(N_{ES} T)$       & $\leq T{+}2$          & 8--12\%  & 0--5  & \checkmark \\
F & Peak-period cover    & $O(|\mathcal{A}|)$  & $\leq N^a$            & 3--5\%   & 0 only& \checkmark \\
D & DC direction         & $O(N_{DC} T)$       & $\leq N^{recv}{+}1$   & 5--10\%  & 0--5  & --- \\
H & Interprovincial sched.& $O(|\mathcal{A}|T)$& $\leq N^a{+}N_{ES}$   & 5--8\%   & 0 only& --- \\
I & PFR--storage         & $O(T)$              & $\leq N{+}N_{ES}$     & 3--5\%   & 0--5  & --- \\
\bottomrule
\end{longtable}

The \checkmark{} column indicates families currently implemented as
pre-formulation rows in \texttt{market\_simulation.cpp}.
Families~D, H, and~I require model elements absent from the current
formulation (DC direction binaries, provincial export schedules, and PFR
requirement data, respectively) and are reserved for future extensions.

Among the implemented families, B, C, and F include an additional
feasibility guard: the cut is skipped for any area whose total
$\sum_{i} P^{\max}_i$ is less than the residual requirement $r_{a,h}$.
In a multi-area system with unconstrained inter-area imports the area
could satisfy demand through imports, so requiring $\sum P^{\max}_i
\geq r_{a,h}$ at the LP level would produce an infeasible LP row.

\paragraph{Relation to generic families.}
Family~G is a strict extension of generic Family~2 (startup clique),
enlarging the window from $T^U_i$ to $T^U_i + T^D_i$.
Family~B incorporates the reserve dimension absent from generic Family~5
and additionally subtracts wind, solar, hydro, and storage non-thermal
maximum output from the requirement.
Family~F is a peak-period variant of B (applied to the single highest-demand
commitment period per area) and replaces the earlier conservative
``hydro-thermal only'' formulation that incorrectly excluded renewable credit.
Family~E complements generic Family~9 by additionally exploiting the
explicit cycle-count parameter $N^{\mathrm{cycle}}_e$.
Families~A and the corrected~F have no direct counterparts in the generic set.

%===================================================================
\subsection{Numerical Performance: Pre-Formulation Cuts on Five-Province SCUC}
\label{sec:market-cuts-perf}
%===================================================================

The cutting-plane implementation was validated on the five-province
day-ahead SCUC benchmark in
\texttt{tests/test\_market\_simulation.cpp}
(\texttt{[market][cuts][benchmark]}).
The test solves the same instance twice---without and with all implemented
pre-formulation cut families (\S\ref{sec:milp-tightening}: 6, 11 and
\S\ref{sec:scuc-market-cuts}: A, B, C, E, F, G)---using Gurobi with default
settings.

\begin{table}[h]
  \centering
  \caption{Five-province SCUC benchmark: baseline vs.\ with pre-formulation
           cuts (Gurobi; 8 commit periods; 96-period 15\,min dispatch).}
  \label{tab:cuts-perf}
  \begin{tabular}{lrr}
    \toprule
    \textbf{Metric} & \textbf{Baseline} & \textbf{With cuts} \\
    \midrule
    Converged                  & yes         & yes         \\
    Objective (\$/interval)    & 321\,921.71 & 321\,921.71 \\
    Objective deviation        & ---         & 0.000\%     \\
    Solve time (s)             & 0.619       & 0.270       \\
    Solve-time change          & ---         & $-56.3\%$   \\
    Final MIP gap              & 0.000\%     & 0.000\%     \\
    Pre-formulation cut rows   & 0           & 1\,869      \\
    \bottomrule
  \end{tabular}
\end{table}

The 1\,869 cut rows are dominated by Family~A (one row per segment per unit
per dispatch period) and Family~G (one clique row per generator per
commitment window).  Families~B and~F each contribute one row per
balancing area per relevant commit period; Families~E and~11 contribute
$2N_{ES}$ rows; Family~C contributes rows only when monitored sections
are present.

The 56\% solve-time reduction is consistent with the expected root-LP
tightening effect: a tighter LP reduces both the initial MIP gap and the
number of branch-and-bound nodes required for proof.
The zero objective deviation confirms that every added row is LP-valid
and does not cut off any feasible integer solution.

\paragraph{Implementation notes.}
Families~1, 2, 3, and~5 (\S\ref{sec:gen-cuts}) are \emph{not} added as
pre-formulation rows.  In the LP relaxation, startup and shutdown
variables $v_{i,t'}$ and $w_{i,t'}$ are fractional; the ramp-cut
right-hand sides can consequently become negative while the lower bound
$P^{\min}_i u_{i,t} > 0$ is simultaneously active, making the LP
infeasible.  These families must be applied as B\&B separation cuts at
nodes where $v_{i,t'} \in \{0,1\}$ is already determined.

Families~B, C, and~F include an inter-area import guard: a cut is
skipped for any area whose committed thermal capacity
$\sum_{i \in \mathcal{G}^a} P^{\max}_i$ falls below the residual
requirement $r_{a,h}$.  Without the guard, areas that rely on
unconstrained imports would receive infeasible LP rows.
